\section{Related work}
\label{sec:related}

\subsection{Playback-conditioned transcript and history truncation}

The closest precedents at the interaction-policy level revise conversation state according to playback progress. OpenAI removes the unplayed audio and transcript suffix at a client-supplied audio boundary \citep{openaiRealtimeConversations}. Azure similarly updates the response and session context after an interruption \citep{microsoftVoiceLiveAutoTruncation2026}, while LiveKit exposes framework-level semantics for truncating interrupted speech \citep{livekitEventsErrorHandling}. These systems establish that an assistant message need not be retained according to generation completion.

Their public interfaces do not establish their inference internals. The reviewed OpenAI and Azure documentation does not disclose whether the services use a cascaded architecture, how assistant token spans are represented, or whether KV state is cropped in place. LiveKit exposes message and playback behavior but not joint recovery of transformer K/V tensors, masks, token ledgers, positions, and role/EOT state. Our distinction is therefore one of inspectable mechanism and evidence, not an assertion that these systems lack an undisclosed implementation. Their playback quantities are also software-level values or estimates, not direct measurements of device presentation or acoustic reception.

\subsection{Streaming dialogue, anticipatory computation, and interruption control}

Predictive ASR uses partial recognition to forecast completed input and prefetch a response that can be used if the final recognition confirms the prediction \citep{schwarz2023PersonalizedPredictiveASR}. LTS-VoiceAgent combines semantic triggering with incremental reasoning \citep{zou2026LTSVoiceAgent}, and RelayS2S uses response-level candidate prefixes with verification and continuation \citep{mai2026RelayS2S}. These studies establish compute-before-commit as a relevant design pattern for speech interaction. Our supporting candidate study shares that pattern but measures generation before synchronous oracle acceptance. Its candidate-selection event is not equivalent to production availability because cache update, delivery authorization, consumer admission, TTS, and audio buffering may follow it.

Interruption-control research primarily determines when output should stop. FireRedChat combines streaming or personalized voice-activity detection with interruption control in cascaded and semi-cascaded configurations \citep{chen2025FireRedChat}. In contrast, joint prefix-state repair starts after an interruption decision and determines what part of the assistant turn remains committed. Detection and repair are complementary, but successful detection does not imply a valid next-turn cache and dialogue history.

Full-duplex speech models provide another reference. Moshi models concurrent speech and text streams to support overlapping interaction \citep{defossez2024Moshi}. Tighter stream coordination can avoid some boundaries introduced by independently deployed ASR, LLM, and TTS components, but model-stream position, device presentation, and acoustic reception remain distinct in the presence of queues and buffers. Comparisons must therefore name the observed boundary instead of inferring delivery from an architectural label.

Related studies connect streaming speech processing, conversational timing, and contextual inference. Morrone et al. integrated speech separation and voice-activity detection for low-latency diarization \citep{morrone2024LowLatencyDiarization}; Edwards measured how telecommunications delay changes speaker-transition timing \citep{edwards2025TelecommunicationsLatency}; and Zheng et al. combined conversation history and hotwords for contextual prompting in conversational ASR \citep{zheng2026MultilevelContextualPrompting}. These works address speech processing, latency, and conversational context, respectively. None of them addresses playback-conditioned transformer-state rollback.

\subsection{KV-cache management and cross-turn state}

Autoregressive transformers retain key and value tensors for prior tokens to avoid recomputing the full prefix. Transformers exposes this state through cache abstractions, including \texttt{DynamicCache.crop} \citep{huggingFace2025TransformersCacheUtils,huggingFace2025KVCacheStrategies}. The primitive makes in-place truncation possible but neither selects the conversational boundary nor repairs non-cache state.

Serving research addresses adjacent objectives. PagedAttention organizes KV memory into non-contiguous blocks with demand allocation and copy-on-write \citep{kwon2023PagedAttention}. SGLang's RadixAttention manages reusable token prefixes in a radix tree \citep{zheng2024SGLang}. These methods improve memory use, throughput, or shared-prefix reuse rather than select an interrupted assistant prefix from software playback progress. IntentKV prunes cross-turn state according to text-agent intent \citep{li2026IntentKV}, and Speculative Interaction Agents coordinate provisional results around asynchronous tool calls \citep{hooper2026SpeculativeInteractionAgents}. Their control signals are text intent or tool state, not TTS-fragment ownership and software-consumed samples.

The technical distinction of the present work is consequently not a new cache representation or crop operator. It couples an external software progress signal to a legal conversational prefix, then repairs all representations of that prefix in one transition. Its validation is equally specific: direct slicing of the same pre-crop snapshot checks whether the requested K/V prefix was retained; a wrong-length control checks whether the gate can reject a boundary error; matched recovery checks whether equal retained states remain equal under identical subsequent token chunks and operations. These gates do not establish equivalence to independently recomputed prefixes, semantic correctness of subsequent responses, or portability across inference engines.

A targeted public-source scan was frozen on 3 September 2026. It covered cross-publisher indexes, arXiv, ACL Anthology, ISCA Archive, ACM Digital Library, IEEE Xplore, NeurIPS proceedings, DOI/Crossref metadata, and first-party product and repository sources. Queries combined cascaded speech dialogue, barge-in, anticipatory response generation, cache rollback, and playback-aware history. Within the reported and accessible sources, no publicly inspectable cascaded implementation was identified that disclosed the complete cursor--fragment--token--cache--role path together with reproducible direct-integrity and recovery evidence. Some sources were inaccessible because of automated-access restrictions, and the search was neither a systematic review nor a patent search. This is a bounded non-identification result, not a global-first claim.
